\chapter{Intravenous (IV) Therapy}

\section*{Overview}
IV therapy delivers fluids, nutrients, or medications directly into a person's vein, allowing immediate absorption and rapid distribution throughout the circulatory system.

\section*{Alternative Names}
IV fluid administration, Intravenous infusion

\section*{Principle}
A cannula placed in a vein delivers fluids, electrolytes, medications, blood products, or nutrition directly into the bloodstream. Absorption is immediate and bypasses the digestive tract, allowing rapid effect and continuous administration.
\section*{History}
Sir Christopher Wren performed the first IV injection in 1656 using a quill and bladder. The development of plastic cannulas in the 1950s made IV therapy safe and routine.

\section*{Why This Test or Procedure Is Ordered}
Ordered to treat dehydration, administer antibiotics, chemotherapy, blood products, or emergency medications that cannot be given orally.

\section*{Is This a Primary or Secondary Test or Procedure}
Primary method for acute fluid resuscitation, electrolyte correction, and urgent drug delivery.

\section*{How This Test or Procedure Is Performed}
A nurse inserts a small plastic catheter (IV line) into a vein, typically in your hand or arm, using a needle which is then removed. Fluids or medications are connected to the catheter.

\section*{What Preparation Is Needed}
No special preparation. Tell the nurse if you are allergic to latex, adhesives, or antiseptics like chlorhexidine.

\section*{Contraindications and Precautions}
Avoid IV insertion in a limb with an active infection, lymphedema, or an AV fistula.

\section*{Risks}
Infiltration (fluid leaking into surrounding tissue), phlebitis (vein inflammation), infection at the insertion site, or fluid overload.

\section*{What Happens After the Test or Procedure}
The catheter is removed, and pressure is applied to the site to prevent bleeding. Monitor for redness or swelling.

\section*{Understanding Your Results}
Successful therapy is indicated by the resolution of symptoms (e.g., rehydration, pain control, or target drug levels).

\section*{Normal Results}
Successful catheter insertion with patent flow; no infiltration, pain, swelling, or redness at the site.

\section*{Follow-up Tests and Procedures}
Monitoring vital signs, fluid balance, and clinical improvement.

\section*{References}
\begin{enumerate}
  \item MedlinePlus. Intravenous (IV) Therapy. U.S. National Library of Medicine; 2025.
  \item Current Medical Diagnosis and Treatment. McGraw-Hill; 2025.
\end{enumerate}

\clearpage
